% Using InitechBase Chapter 5, entries ERASE .. RETURN (AN-04).

% ERASE -- os/samir/cmd/mutate.c d_file (Using III+ U5-121; a file of an
% open database is refused, 3; absent, 1).  Lane samirdml 14a8fc75.
% Gates: test-samir-dml, test-samir-dml-emu.
\begin{XbaseEntry}{ERASE}{command}
Deletes a file from the diskette.
\Syntax
\SyntaxLine{\kw{ERASE} \xp{file}}
\Usage
\cmd{ERASE} deletes the file named.  Give the file name with its extension;
none is assumed.  The file cannot be recovered.

A database that is open in a work area cannot be erased:
\cmd{ERASE} reports \xmsg{3}{File is already open.}  Close it with \cmd{USE}
or \cmd{CLOSE} first.  If the file is not on the diskette, \cmd{ERASE}
reports \xmsg{1}{File does not exist.}
\SeeAlso{RENAME, TYPE}
\end{XbaseEntry}

% FIND -- os/samir/cmd/query.c q_find (unquoted literal through ndx_seek on
% the master index; SET EXACT; #26 without an index).  Gate: test-ndx-seek,
% test-samir-query.
\begin{XbaseEntry}{FIND}{command}
Finds the first record whose index key begins with a literal character
string.
\Syntax
\SyntaxLine{\kw{FIND} \xp{literal string}}
\Usage
\cmd{FIND} searches the master index of the database in USE (the first index
named in \cmd{USE}, or the one chosen with SET ORDER) for the characters
typed after the command, without quotation marks.  If a key begins with them,
the record pointer moves to the first such record in index order and
\fn{FOUND()} is true.  If none does, the pointer moves to the end of the file:
\fn{EOF()} is true and \fn{FOUND()} false.

While SET EXACT is OFF, a key that begins with the string is found; with SET
EXACT ON, the key must equal the string, apart from trailing blanks.  Capital
and small letters are different.

\cmd{FIND} needs an open index.  Without one it displays
\xmsg{26}{Database is not indexed.}
\Tips
\cmd{SEEK} does the same with an expression, so it can search for the
contents of a memory variable or for a number.
\Examples
\DotExample{examples/c-find}
A string that begins no key leaves the pointer at the end of the file:
\DotExample{examples/c-find-2}
\SeeAlso{SEEK, LOCATE, SET EXACT, SET ORDER, USE, FOUND()}
\end{XbaseEntry}

% GO -- os/samir/cmd/query.c q_go (TOP / BOTTOM in master-index order, a
% record number in physical order; #5 out of range).  Gate: test-interp-nav.
\begin{XbaseEntry}{GO}{command}
\XAlso{GOTO}{GOTO}%
Moves the record pointer to a record.
\Syntax
\SyntaxLine{\kw{GO}/\kw{GOTO} \xp{expN}}
\par or\par
\SyntaxLine{\kw{GO}/\kw{GOTO} \kw{TOP}/\kw{BOTTOM}}
\Usage
\cmd{GO} \xp{expN} makes record number \xp{expN} the current record.  The
record number is the position of the record in the database, whatever index
is open.  A number below 1 or above the number of records is refused with
\xmsg{5}{Record is out of range.}

\cmd{GO TOP} moves to the first record and \cmd{GO BOTTOM} to the last.  With
an index open, first and last mean in the order of the master index; without
one, record 1 and the last record.  While a filter is set (SET FILTER),
\cmd{GO TOP} and \cmd{GO BOTTOM} move to the first and last record that meets
it.

\cmd{GOTO} is another name for \cmd{GO}.
\Examples
\DotExample{examples/c-go}
With an index open, the top is the first name in alphabetical order:
\DotExample{examples/c-go-2}
\SeeAlso{SKIP, RECNO(), EOF(), BOF()}
\end{XbaseEntry}

% IF -- os/samir/cmd/flow.c LK_IF/LK_ELSE/LK_ENDIF (guard must be Logical,
% #37; program lines only).  Gate: test-interp-flow.
\begin{XbaseEntry}{IF}{command}
\XAlso{ELSE}{ELSE}\XAlso{ENDIF}{ENDIF}%
Carries out one of two groups of commands in a program, according to a
condition.
\Syntax
\SyntaxLine{\kw{IF} \xp{condition}}
\SyntaxLine{\ \ \xp{commands}}
\SyntaxLine{\opt{\kw{ELSE}}}
\SyntaxLine{\ \ \opt{\xp{commands}}}
\SyntaxLine{\kw{ENDIF}}
\Usage
When the condition is true, the commands between \cmd{IF} and \cmd{ELSE} (or
\cmd{ENDIF}) are carried out.  When it is false, the commands between
\cmd{ELSE} and \cmd{ENDIF} are carried out, if there is an \cmd{ELSE}.  The
program then goes on after \cmd{ENDIF}.

The condition must be a logical expression: a field of type Logical, a
comparison, or a logical function such as \fn{EOF()}.  InitechBase does not
treat a number or a string as true or false; a condition of another type stops
the program with \xmsg{37}{Not a Logical expression.}

\cmd{IF} is used only in programs.  \cmd{IF} structures may be nested inside one
another and inside \cmd{DO WHILE} and \cmd{DO CASE}; each \cmd{IF} needs its own
\cmd{ENDIF}.
\Examples
\file{PAYTOTAL.PRG}, listed under \cmd{DO WHILE}, uses \cmd{IF} twice: to skip
the deleted records and to add each salary to one of two totals:
\DotExample{examples/c-if}
\SeeAlso{DO CASE, DO WHILE, IIF()}
\end{XbaseEntry}

% INDEX -- os/samir/cmd/mutate.c d_index (ON <key> TO <file>, default .NDX;
% key Character, Numeric or Date, 1..220 bytes of text; the new index is
% attached and becomes the master; SAFETY ON refuses an existing file, 7).
% Lane samirdml 14a8fc75.  Gates: test-samir-dml, test-samir-dml-emu.
\begin{XbaseEntry}{INDEX}{command}
Creates an index file for the database in USE.
\Syntax
\SyntaxLine{\kw{INDEX ON} \xp{key expression} \kw{TO} \xp{index file}}
\Usage
\cmd{INDEX} builds an index file of the database in the selected work area,
ordered by the key expression.  The key must be of type Character, Numeric
or Date.  If no extension is given, \file{.NDX} is assumed.  The new index is
opened at once, with the index files already open, as the master index; the
record pointer goes to the first record in key order.  Up to seven index
files may be open with one database.

While SET SAFETY is ON, an index file that already exists is not replaced:
\cmd{INDEX} reports \xmsg{7}{File already exists.}
\SeeAlso{REINDEX, SET INDEX, SET ORDER, USE}
\end{XbaseEntry}

% INPUT -- os/samir/cmd/proc.c p_input (prompt, one line, EVALUATED as an
% expression; the memvar takes its type).  Gate: test-interp-proc.
\begin{XbaseEntry}{INPUT}{command}
Displays a prompt, reads an expression typed at the keyboard and stores its
value in a memory variable.
\Syntax
\SyntaxLine{\kw{INPUT} \opt{\xp{expC}} \kw{TO} \xp{memvar}}
\Usage
\cmd{INPUT} displays the prompt \xp{expC}, if one is given, and waits for a
line ending with \key{Enter}.  The line is evaluated as an expression, and its
value is stored in \xp{memvar}, which takes the type of the value.  A number
typed in answer gives a numeric memory variable.  A character string must be
typed in quotation marks; a date is typed as a call of \fn{CTOD()}; a logical
value as \texttt{.T.} or \texttt{.F.}.

Because the answer is an expression, it may also be a calculation, a field
name or a memory variable.  The answer must be a valid InitechBase
expression.
\Tips
Use \cmd{ACCEPT} for character answers: the user need not type quotation
marks.
\Examples
\DotExample{examples/c-input}
\SeeAlso{ACCEPT, WAIT, STORE, CTOD()}
\end{XbaseEntry}

% LABEL -- os/samir/cmd/report.c lbl_run (FORM <file>, default .LBL, the
% 1034-byte dBASE III PLUS label form; TO FILE <file>, default .TXT; every
% record in master-index order; 54 for an invalid form).  Lane samirdml
% 14a8fc75.  Gates: test-samir-dml, test-samir-dml-emu.
\begin{XbaseEntry}{LABEL}{command}
Prints labels from a label form file.
\Syntax
\SyntaxLine{\kw{LABEL FORM} \xp{label file} \opt{\kw{TO FILE} \xp{file}}}
\Usage
\cmd{LABEL FORM} prints one label for every record of the database in USE,
in the order of the master index, laid out by the label form file named.
If no extension is given, \file{.LBL} is assumed.  The label form gives the
height and width of a label, the margin, the space between labels, the
number of labels across, and the expressions printed on each line of the
label.

\kw{TO FILE} writes the labels to a text file instead of the screen;
\file{.TXT} is assumed.

InitechBase prints label forms made with dBASE III PLUS; it does not create
them.  A form that is not a valid label form is reported as
\xmsg{54}{Label file invalid.}
\SeeAlso{REPORT FORM}
\end{XbaseEntry}

% LIST -- os/samir/cmd/query.c q_list (default scope ALL with an implicit
% GO TOP; heading per SET HEADING; deleted '*'; SET DELETED, SET FILTER).
% Gates: test-samir-query, test-samir-list-layout, test-interp-list.
\begin{XbaseEntry}{LIST}{command}
Lists the records of the database in USE.
\Syntax
\SyntaxLine{\kw{LIST} \opt{\xp{scope}} \opt{\xp{expression list}}
  \opt{\kw{FOR} \xp{condition}} \opt{\kw{WHILE} \xp{condition}}
  \opt{\kw{OFF}}}
\Defaults
The scope is \kw{ALL}: \cmd{LIST} begins at the first record.  Without an
expression list, every field is listed.
\Usage
\cmd{LIST} displays one line for each record of the scope that meets the
\kw{FOR} and \kw{WHILE} conditions.  Each line begins with the record number,
then the value of each field or expression.  The first line is a heading:
\texttt{Record\#} and the field names, or the expressions as typed; SET
HEADING OFF leaves it out.  A record marked for deletion has an asterisk in
front of its number.  While SET DELETED is ON, such records are not listed;
while a filter is set, only the records that meet it are listed.

The columns are as wide as the fields or the values, or as the headings when
these are wider.  Character fields are shown at their full width, numeric
fields with their decimal places, dates in the form set by SET DATE, logical
fields as \texttt{T} or \texttt{F}, and memo fields with their text.  The
screen holds 80 characters: a longer line continues on the next.

\kw{FOR} tests every record of the scope.  \kw{WHILE} stops at the first record
for which its condition is false, and leaves the record pointer on it; without
a scope, a \cmd{LIST} with \kw{WHILE} also begins at the first record.  The
scope must come before the expression list.  After \cmd{LIST} of the whole
file the record pointer is at the end of the file; after \kw{NEXT} \xp{n} it is
on the last record listed.
\Options
\kw{OFF} leaves out the record numbers.  In this version, \kw{OFF} lists
every field of each record and takes no expression list.
\Examples
\DotExample{examples/c-list}
Expressions as well as fields may be listed:
\DotExample{examples/c-list-2}
\SeeAlso{DISPLAY, ?, SET DELETED, SET FILTER, SET HEADING, REPORT FORM}
\end{XbaseEntry}

% LOCATE -- os/samir/cmd/query.c q_locate (default scope ALL with implicit
% GO TOP; stores the FOR for CONTINUE; sets FOUND/EOF).  Gate:
% test-samir-query.
\begin{XbaseEntry}{LOCATE}{command}
Moves to the first record that meets a condition.
\Syntax
\SyntaxLine{\kw{LOCATE} \opt{\xp{scope}} \kw{FOR} \xp{condition}
  \opt{\kw{WHILE} \xp{condition}}}
\Defaults
The scope is \kw{ALL}: the search begins at the first record.
\Usage
\cmd{LOCATE} tests the records of the scope, in order, against the \kw{FOR}
condition, and stops at the first record that meets it.  That record becomes
the current record and \fn{FOUND()} is true.  If no record meets the
condition, the record pointer moves to the end of the scope: for the default
scope, \fn{EOF()} is true.  \fn{FOUND()} is then false.

\cmd{LOCATE} reads the records one after another and needs no index.  The
condition may be any logical expression.  \cmd{CONTINUE} goes on to the next
record that meets the same condition.
\Tips
On a large indexed database, \cmd{SEEK} finds a key faster than \cmd{LOCATE}.
\Examples
\DotExample{examples/c-locate}
\SeeAlso{CONTINUE, SEEK, FIND, SET FILTER, FOUND()}
\end{XbaseEntry}

% NOTE / * / && -- os/samir/cmd/flow.c trim_line ("&&" strips to the end of
% the line; a line beginning "*" or "NOTE" is a comment).  Gate:
% test-interp-flow.
\begin{XbaseEntry}{NOTE}{command}
\XAlso{*}{*}\XAlso{&&}{\&\&}%
Marks a comment in a program.
\Syntax
\SyntaxLine{\kw{NOTE} \xp{text}}
\par or\par
\SyntaxLine{\kw{*} \xp{text}}
\par or\par
\SyntaxLine{\xp{command} \kw{\&\&} \xp{text}}
\Usage
A line that begins with \cmd{NOTE} or with an asterisk is a comment:
InitechBase does not carry it out.  Two ampersands after a command begin a
comment that runs to the end of the line; the command before them is carried
out.

Comments explain a program to the people who read it.  Every program on the
Sample Programs and Utilities diskette begins with comment lines that give its
name, its purpose and the database it works on.
\Examples
\DotExample{examples/c-note}
\SeeAlso{DO, PROCEDURE}
\end{XbaseEntry}

% ON ERROR -- os/samir/cmd/proc.c p_on "ERROR" (store / clear the trap) and
% proc_run (the trap command runs once on a run-time error; the message is
% then displayed).  Gate: test-interp-proc.
\begin{XbaseEntry}{ON ERROR}{compound}
\XAlso{ON}{ON}%
Names a command to be carried out when an error occurs.
\Syntax
\SyntaxLine{\kw{ON ERROR} \opt{\xp{command}}}
\Usage
After \cmd{ON ERROR} \xp{command}, whenever a command or a program stops with
an error, InitechBase carries out \xp{command}, and then displays the error
message as usual.  The function \fn{ERROR()} gives the number of the error
while the command runs, so the command can record or report it.  The command is
often \cmd{DO} of a procedure that handles errors.

\cmd{ON ERROR} with no command removes the trap.  A new \cmd{ON ERROR} replaces
the one before.
\Examples
\DotExample{examples/c-onerror}
\SeeAlso{ERROR(), ON ESCAPE, ON KEY, RETURN}
\end{XbaseEntry}

% ON ESCAPE, ON KEY -- os/samir/cmd/proc.c p_on "ESCAPE"/"KEY": parsed and
% accepted, the trap is not installed.
\begin{XbaseEntry}[the ON ESCAPE and ON KEY traps are accepted but not
  carried out by this version (proc.c)]{ON ESCAPE}{compound}
Name a command to be carried out when \key{Esc}, or any key, is pressed while
a program runs.
\Syntax
\SyntaxLine{\kw{ON ESCAPE} \opt{\xp{command}}}
\par or\par
\SyntaxLine{\kw{ON KEY} \opt{\xp{command}}}
\Usage
InitechBase accepts both commands, with or without a command to carry out,
and displays no message.  The command named is not carried out when a key is
pressed.
\SeeAlso{ON ERROR, WAIT}
\end{XbaseEntry}

% PACK -- os/samir/cmd/mutate.c m_pack (dbf_pack removes the 0x2A records and
% renumbers; every open index rebuilt, initech-x87g).  Gates:
% test-interp-replace, test-ndx-pack.
\begin{XbaseEntry}{PACK}{command}
Removes the records marked for deletion from the database in USE.
\Syntax
\SyntaxLine{\kw{PACK}}
\Usage
\cmd{PACK} removes every record that \cmd{DELETE} has marked from the database
in the selected work area, and renumbers the records that remain, in their
order.  The database becomes smaller, and \fn{RECCOUNT()} gives the new
number of records.  Every index file open with the database is rebuilt for the
new record numbers.

The records removed cannot be recalled.
\Tips
Before \cmd{PACK}, list the marked records with \cmd{LIST FOR DELETED()} to make
sure no record is marked by mistake.
\Examples
\DotExample{examples/c-pack}
\SeeAlso{DELETE, RECALL, ZAP, DELETED()}
\end{XbaseEntry}

% PARAMETERS -- os/samir/cmd/proc.c p_parameters (binds DO ... WITH values
% by position into PRIVATE memvars; #93 without arguments, #94 count).
% Gate: test-interp-proc.
\begin{XbaseEntry}{PARAMETERS}{command}
\XAlso{PARAMETER}{PARAMETER}%
Receives the values passed by \cmd{DO} \ldots{} \kw{WITH} in memory variables.
\Syntax
\SyntaxLine{\kw{PARAMETERS} \xp{memvar list}}
\Usage
\cmd{PARAMETERS} is the first command of a procedure that is called with
\kw{WITH}.  It creates one private memory variable for each name of the list
and gives it the value in the same position of the \kw{WITH} list: the first
name receives the first value, and so on.  The memory variables are released
when the procedure returns.

Give \cmd{PARAMETERS} one name for each value the \cmd{DO} passes.  A procedure
called with values but without \cmd{PARAMETERS} stops at its \cmd{RETURN} with
\xmsg{93}{No PARAMETER statement found.}

\cmd{PARAMETER} is another name for \cmd{PARAMETERS}.
\Examples
See \cmd{DO}: the procedure \texttt{NEWSAL} of \file{RAISE.PRG} receives a salary
and a percentage.
\SeeAlso{DO, PROCEDURE, PRIVATE, RETURN}
\end{XbaseEntry}

% PRIVATE -- os/samir/cmd/proc.c p_private -> xb_interp_declare_private
% (hides the caller's variable until the routine returns).  Gate:
% test-interp-proc.
\begin{XbaseEntry}{PRIVATE}{command}
Creates memory variables that belong to the procedure running, hiding any of
the same names.
\Syntax
\SyntaxLine{\kw{PRIVATE} \xp{memvar list}}
\Usage
\cmd{PRIVATE} gives the procedure in which it is used its own memory variables
of the names in the list.  A memory variable of the same name created by the
program that called the procedure, or declared PUBLIC, is hidden while the
procedure runs and keeps its value: when the procedure returns, the private
variables are released and the hidden ones are visible again.

A private memory variable has no value until a value is stored in it.

A memory variable that a procedure creates without \cmd{PRIVATE}, by storing a
value in a name that does not yet exist, is private to that procedure as well.
\Examples
\file{TALLY.PRG} keeps its count apart from the count of its procedure:
\ProgramListing{TALLY.PRG}
\DotExample{examples/c-private}
\SeeAlso{PUBLIC, RELEASE, STORE, PARAMETERS}
\end{XbaseEntry}

% PROCEDURE -- os/samir/cmd/proc.c (PROCEDURE lines end the main program;
% DO <name> runs the block up to RETURN).  Gate: test-interp-proc.
\begin{XbaseEntry}{PROCEDURE}{command}
Begins a procedure in a program file.
\Syntax
\SyntaxLine{\kw{PROCEDURE} \xp{name}}
\Usage
A program file may hold, after its main program, any number of procedures.
Each begins with a \cmd{PROCEDURE} line that names it, and ends with
\cmd{RETURN}.  The main program ends at the first \cmd{PROCEDURE} line: when a
program is run, it stops there, and the procedures run only when
\cmd{DO} calls them by name.

A procedure name has up to 10 letters and digits and begins with a letter.
\Examples
See \cmd{DO} and \cmd{PRIVATE}.
\SeeAlso{DO, PARAMETERS, RETURN}
\end{XbaseEntry}

% PUBLIC -- os/samir/cmd/proc.c p_public -> xb_interp_declare_public
% (level-0 variable that survives RETURN).  Gate: test-interp-proc.
\begin{XbaseEntry}{PUBLIC}{command}
Creates memory variables that every program and procedure can use and that
remain when they return.
\Syntax
\SyntaxLine{\kw{PUBLIC} \xp{memvar list}}
\Usage
A public memory variable is not released when the procedure or program that
created it returns.  It can be used by every procedure, by the programs run
afterwards and at the dot prompt, until it is released with \cmd{RELEASE}.

Declare a memory variable \cmd{PUBLIC} before you store a value in it.  If a
memory variable of the name is already visible, \cmd{PUBLIC} leaves it as it
is.  A public memory variable has no value until one is stored in it.
\Examples
The procedure \texttt{NEWSAL} of \file{RAISE.PRG} (listed under \cmd{DO})
declares its result \cmd{PUBLIC}, so it is still there at the dot prompt:
\DotExample{examples/c-public}
\SeeAlso{PRIVATE, RELEASE, STORE}
\end{XbaseEntry}

% QUIT -- os/samir/samir_main.c samir_repl (QUIT and EXIT end the loop;
% wa_close_all closes every database; back to InitechDOS).  Gate:
% test-samir-repl, test-samir-boot.
\begin{XbaseEntry}{QUIT}{command}
Closes all files and leaves InitechBase.
\Syntax
\SyntaxLine{\kw{QUIT}}
\Usage
\cmd{QUIT} closes every open database with its index files and returns to
InitechDOS, which displays its prompt.  \cmd{EXIT}, typed at the dot prompt,
does the same.

Always leave InitechBase with \cmd{QUIT} before you turn off or restart the
computer.
\Examples
\DotExample{examples/c-quit}
\SeeAlso{CLOSE, USE}
\end{XbaseEntry}

% RECALL -- os/samir/cmd/mutate.c m_recall (default the current record; FOR
% without a scope = ALL).  Gate: test-interp-replace.
\begin{XbaseEntry}{RECALL}{command}
Removes the deletion mark from records.
\Syntax
\SyntaxLine{\kw{RECALL} \opt{\xp{scope}} \opt{\kw{FOR} \xp{condition}}
  \opt{\kw{WHILE} \xp{condition}}}
\Defaults
Without a scope or a condition, \cmd{RECALL} works on the current record
only.  With a \kw{FOR} condition and no scope, the scope is \kw{ALL}.
\Usage
\cmd{RECALL} makes records marked by \cmd{DELETE} ordinary records again.  A
record removed by \cmd{PACK} or \cmd{ZAP} cannot be recalled.
\Examples
Record 13 of \file{STAFF.DBF} is marked for deletion on the diskette:
\DotExample{examples/c-recall}
\SeeAlso{DELETE, PACK, DELETED()}
\end{XbaseEntry}

% REINDEX -- os/samir/cmd/mutate.c d_index (rebuild every open index of the
% selected work area from its own key; 114 on failure).  Lane samirdml
% 14a8fc75.  Gates: test-samir-dml, test-samir-dml-emu.
\begin{XbaseEntry}{REINDEX}{command}
Rebuilds the open index files of the database in USE.
\Syntax
\SyntaxLine{\kw{REINDEX}}
\Usage
\cmd{REINDEX} rebuilds every index file open in the selected work area from
its own key expression.  Use it when an index file has been damaged, or
when the database was changed while the index file was not open.
\SeeAlso{INDEX, PACK, USE}
\end{XbaseEntry}

% RELEASE -- os/samir/cmd/flow.c "RELEASE" (a name list, or ALL).  Gate:
% test-interp-flow.
\begin{XbaseEntry}{RELEASE}{command}
Releases memory variables.
\Syntax
\SyntaxLine{\kw{RELEASE} \xp{memvar list}}
\par or\par
\SyntaxLine{\kw{RELEASE ALL}}
\Usage
\cmd{RELEASE} removes the memory variables named.  \cmd{RELEASE ALL} removes all
of them.  A memory variable released no longer exists; a reference to it
gives \xmsg{12}{Variable not found.}
\Examples
\DotExample{examples/c-release}
\DotExample{examples/c-release-2}
\SeeAlso{STORE, PUBLIC, PRIVATE}
\end{XbaseEntry}

% RENAME -- os/samir/cmd/mutate.c d_file (<file> TO <file>; open database
% refused, 3; SAFETY ON refuses an existing target, 7).  Lane samirdml
% 14a8fc75.  Gates: test-samir-dml, test-samir-dml-emu.
\begin{XbaseEntry}{RENAME}{command}
Changes the name of a file.
\Syntax
\SyntaxLine{\kw{RENAME} \xp{file} \kw{TO} \xp{new file}}
\Usage
\cmd{RENAME} gives the file a new name.  Give both names with their
extensions; none is assumed.  A database that is open cannot be renamed
(\xmsg{3}{File is already open.}), and while SET SAFETY is ON a file of the
new name is not replaced (\xmsg{7}{File already exists.}).

Rename a database's memo file (\file{.DBT}) with it.
\SeeAlso{ERASE, TYPE}
\end{XbaseEntry}

% REPLACE -- os/samir/cmd/mutate.c m_replace (default the current record;
% field WITH expr pairs; xbase_coercion.json assignment rules via
% dbf_replace; #9 cross-type; numeric overflow '*'; indexes updated on the
% current-record form).  Gates: test-interp-replace, test-dbf-mutate.
\begin{XbaseEntry}{REPLACE}{command}
Changes the contents of fields.
\Syntax
\SyntaxLine{\kw{REPLACE} \opt{\xp{scope}} \xp{field} \kw{WITH} \xp{exp}
  \opt{\kw{,} \xp{field} \kw{WITH} \xp{exp} \rep}
  \opt{\kw{FOR} \xp{condition}} \opt{\kw{WHILE} \xp{condition}}}
\Defaults
Without a scope or a condition, \cmd{REPLACE} changes the current record
only.  With a \kw{FOR} condition and no scope, the scope is \kw{ALL}.
\Usage
For each record of the scope that meets the conditions, \cmd{REPLACE}
evaluates each expression and stores its value in the field before
\kw{WITH}.  The expressions are evaluated for the record being changed, so
\texttt{SALARY WITH SALARY * 1.03} raises each salary by three per cent.  The
changed records are written to the disk at once.

The value must be of the type of the field: a character value into a
character field, a number into a numeric field, a date into a date field, a
logical value into a logical field.
\draft{a character value replaced into a memo field is not stored by this
version, and the text the field held is lost (mutate.c); change memo text with
another program}  Any other combination is
refused with \xmsg{9}{Data type mismatch.} (Appendix~E).  A character value
longer than the field is cut off at the width of the field; a shorter one is
filled with blanks.  A number too large for its field is stored as asterisks.

When the current record is replaced, every index file open with the database
is updated for the new contents of the fields.
\Tips
Convert values of another type first: \fn{STR()} turns a number into
characters, \fn{VAL()} characters into a number, \fn{CTOD()} characters into
a date.
\Examples
\DotExample{examples/c-replace}
A scope and a condition change several records at once:
\DotExample{examples/c-replace-2}
A value of the wrong type is refused:
\DotExample{examples/c-replace-3}
\SeeAlso{APPEND, STORE, STR(), VAL(), CTOD()}
\end{XbaseEntry}

% REPORT FORM -- os/samir/cmd/report.c (U5-177 options; the .FRM of
% frm.md; mint 004 layout; TO PRINT refused #126).  Gates:
% test-samir-report, test-samir-report-emu.
\begin{XbaseEntry}{REPORT FORM}{compound}
\XAlso{REPORT}{REPORT}%
Prints a report of the database in USE, laid out by a report form file.
\Syntax
\SyntaxLine{\kw{REPORT FORM} \xp{file} \opt{\xp{scope}}
  \opt{\kw{FOR} \xp{condition}} \opt{\kw{WHILE} \xp{condition}}
  \opt{\kw{PLAIN}} \opt{\kw{HEADING} \xp{expC}} \opt{\kw{NOEJECT}}
  \opt{\kw{SUMMARY}} \opt{\kw{TO FILE} \xp{file}}}
\Defaults
The scope is \kw{ALL}.  The extension of the report form file is
\file{.FRM}.
\Usage
A report form file describes a report: its title, its columns, the expression
and heading of each column, which columns are totalled, and whether the
records are grouped.  The report forms \file{STAFFRPT.FRM} and
\file{INVRPT.FRM} are on the Sample Programs and Utilities diskette.

\cmd{REPORT FORM} prints a page heading -- the page number, the date and the
title -- then the column headings, then one line for each record of the scope
that meets the conditions, and at the end the total of each totalled column.
In a grouped report a new group begins whenever the group expression changes;
each group has its own heading and subtotals.  Use the report with the records
in the order of the group, through an index.

A column value too long for its column continues on the next line of the
report.  A new page begins when the page is full.
\Options
\kw{PLAIN} leaves out the page number and the date, and prints the title and
the column headings on the first page only.

\kw{HEADING} \xp{expC} prints the value of \xp{expC} after the page number.

\kw{NOEJECT} leaves out the page advance before the first page.

\kw{SUMMARY} prints only the group headings, the subtotals and the total,
without the lines of the records.

\kw{TO FILE} \xp{file} also writes the report to a text file.  Without an
extension, the file is given the extension \file{.TXT}.  The report is
displayed on the screen as well.
\Tips
Print a report on a printer from InitechDOS: write it \kw{TO FILE}, leave
InitechBase and copy the file to the printer.  \kw{TO PRINT} is refused with
\xmsg{126}{Printer is either not connected or turned off.}
\Examples
\DotExample{examples/c-report}
The whole register, without page headings:
\DotExample{examples/c-report-2}
Subtotals only, with a heading:
\DotExample{examples/c-report-3}
\SeeAlso{LIST, SET HEADING, USE}
\end{XbaseEntry}

% RESTORE -- os/samir/cmd/flow.c fl_memfile (FROM <file>, default .MEM;
% ADDITIVE keeps the present variables; the whole file is checked before
% any variable changes, 55).  Lane samirdml 14a8fc75.  Gates:
% test-samir-dml, test-samir-dml-emu.
\begin{XbaseEntry}{RESTORE}{command}
Loads memory variables from a memory file.
\Syntax
\SyntaxLine{\kw{RESTORE FROM} \xp{file} \opt{\kw{ADDITIVE}}}
\Usage
\cmd{RESTORE} reads the memory variables saved with \cmd{SAVE}.  If no
extension is given, \file{.MEM} is assumed.  Without \kw{ADDITIVE}, all
present memory variables are released first; with it, they are kept and the
saved ones are added.

A file that is not a valid memory file is reported as
\xmsg{55}{Memory Variable file is invalid.} and no variable is changed.
\SeeAlso{RELEASE, SAVE}
\end{XbaseEntry}

% RETURN -- os/samir/cmd/proc.c PROC_RC_RETURN (leaves the routine, out of any
% IF / DO WHILE; at the main level ends the program).  Gate:
% test-interp-proc.
\begin{XbaseEntry}{RETURN}{command}
Ends a procedure or a program and returns to the one that called it.
\Syntax
\SyntaxLine{\kw{RETURN}}
\Usage
In a procedure, \cmd{RETURN} ends the procedure at once, even inside an
\cmd{IF} or a \cmd{DO WHILE}, and the calling program goes on with the line
after its \cmd{DO}.  The private memory variables of the procedure are
released.  In the main program of a program file, \cmd{RETURN} ends the program
and InitechBase returns to the dot prompt.

Every procedure ends with \cmd{RETURN}, and so does the main program of a file
that contains procedures.  At the dot prompt, \cmd{RETURN} does nothing.
\Examples
See \cmd{DO}.
\SeeAlso{DO, PROCEDURE, PRIVATE}
\end{XbaseEntry}
